% Figaro Paper M: Behavioral Game Theory of the Two-Mechanism Design
% Psychology and Decisions Science register.
% Audience: behavioral game theorists, decision-science researchers,
% experimental economists, interface-design and incentive-design
% scholars, behavioral-mechanism-design community.
% Preprint, May 2026

\documentclass[11pt,a4paper]{article}

\usepackage[utf8]{inputenc}
\usepackage[T1]{fontenc}
\usepackage{amsmath,amssymb,amsthm}
\usepackage{mathtools}
\usepackage{booktabs}
\usepackage{graphicx}
\usepackage{hyperref}
\usepackage{cleveref}
\usepackage{enumitem}
\usepackage{xcolor}
\usepackage[margin=1in]{geometry}
\usepackage{natbib}
\bibliographystyle{plainnat}

\newtheorem{theorem}{Theorem}[section]
\newtheorem{proposition}[theorem]{Proposition}
\newtheorem{corollary}[theorem]{Corollary}
\theoremstyle{definition}
\newtheorem{definition}[theorem]{Definition}
\newtheorem{example}[theorem]{Example}
\theoremstyle{remark}
\newtheorem{remark}[theorem]{Remark}

\newcommand{\figaro}{\textsc{Figaro}}
\newcommand{\pay}{P}
\newcommand{\cumval}{G}
\DeclareMathOperator{\Coop}{C}
\DeclareMathOperator{\Def}{D}

\title{%
  Behavioral Game Theory of the Two-Mechanism Bonded Commitment:\\
  Loss Aversion, Coordination, Peer Pressure, and Incentive Legibility%
}

\author{
  Alessandro Daliana\thanks{Corresponding author. \texttt{adaliana@figaro.org}}
}

\date{May 2026}

\begin{document}

\maketitle

% ════════════════════════════════════════════════════════════════════
\begin{abstract}
The bonded settlement primitive admits a mechanism-design
analysis as a full-rationality mechanism: cooperation is the
unique profile surviving iterated elimination of weakly dominated
strategies, and the analysis assumes complete information, no
non-pecuniary preferences, and no coordination failure under
pessimistic beliefs. The full-rationality argument is correct on its
own terms; this paper treats the same primitive from a
behavioral-game-theory perspective and asks how the
full-rationality analysis carries over when participants exhibit
the documented departures from full rationality the experimental
literature has catalogued over the past forty years.

The paper develops four claims. First, the primitive's reliance on
weak dominance --- not strict dominance --- is more behaviorally robust
than it might appear, because the relevant cognitive operation is a
single comparison at the margin between cooperation and defection,
amplified by loss aversion at the kink (not at the magnitude). Second,
the Van Huyck--Battalio--Beil coordination-failure result, which
appears to threaten the multi-party reading of the architecture, does
not transfer because the per-player game in the bonded setting is
dominance-solvable in a way the experimental weakest-link games are
not. Third, the externality structure of atomic resolution produces
peer pressure that is both economically real (the mechanism paper's
analytical result) and behaviorally activated (the experimental
literature on social-pressure responses). Fourth, the architecture's
\emph{incentive legibility} --- the property that participants can
verify the equilibrium argument with a single comparison rather than
through iterated reasoning about counterparty rationality --- makes it
unusually well-suited to the cognitive constraints of real
participants.

The paper closes with the empirical questions the architecture
invites: experimental measurement of cooperation rates relative to
controls without bonded backing, behavioral measurement of the
weakest-link pressure activation across process-tree depths, and
field measurement of the cognitive load the architecture imposes
relative to the alternatives it is intended to displace.
\end{abstract}

\medskip
\noindent\textbf{Keywords:}
behavioral game theory, loss aversion, weakest-link coordination,
peer pressure, incentive legibility, mechanism design,
experimental economics, interface cognition

% ════════════════════════════════════════════════════════════════════
\section{Introduction}
\label{sec:intro}

Behavioral game theory \citep{camerer2003behavioral} is the
empirical correction to mechanism design's full-rationality
assumptions. Forty years of experimental work has documented
systematic departures from the rational-choice baseline: loss
aversion \citep{kahneman_tversky_1979,
tversky_kahneman_1992}, framing effects, anchoring, hyperbolic
discounting, fairness preferences \citep{fehr_schmidt_1999,
bolton_ockenfels_2000}, social preferences, and coordination
failures under pessimistic beliefs about co-players
\citep{vanhuyck_battalio_beil_1990}. A mechanism designed under
the rational-choice assumption is robust to behavioral
departures only if the departures preserve the mechanism's
load-bearing equilibrium properties; many do not.

The mechanism-design treatment of the bonded primitive proves
that cooperation is weakly dominant for both parties in the
bonded game, and is the unique profile surviving
iterated elimination of weakly dominated strategies. The proof
assumes complete information, full rationality, no non-pecuniary
preferences, and no coordination failure under pessimistic
beliefs. The result holds within those assumptions. The present
paper asks: how does the result carry over when the assumptions
are weakened in the directions the experimental literature has
identified as empirically relevant?

\paragraph{The four claims.}
The paper develops four claims about the behavioral robustness of
the bonded architecture.

First, the bond's reliance on weak rather than strict dominance
is structurally consequential because the strict-dominance gap is
at the kink between cooperation (recovering the bond) and
defection (forfeiting the bond), not at any magnitude comparison
within either domain. Loss aversion amplifies the kink and is
robust to the magnitude-domain diminishing sensitivity that
prospect theory documents.

Second, the multi-party result is robust to the
\citet{vanhuyck_battalio_beil_1990} coordination-failure pattern
because the per-player game is dominance-solvable in the bonded
setting. The classical experimental result requires individual
non-dominance --- a setting in which pessimistic beliefs about
co-players make defection rational. The bonded setting forecloses
this.

Third, the weakest-link subgame the mechanism paper analyzes is
both economically real and behaviorally activated. The
externality magnitude $\pay_i + 2\cumval_i$ is large; the
experimental literature on peer pressure shows that social-pressure
activation occurs at much lower stakes. The architecture's
peer-enforcement effect should therefore be expected to operate
through both the economic-rationality channel (Paper~A) and the
behavioral peer-pressure channel.

Fourth, the architecture exhibits \emph{incentive legibility}: the
cognitive operation a participant must perform to assess the
equilibrium is a single comparison at the margin. This is the
minimal reasoning task in any strategic setting, and it is well
within the documented capabilities of unsophisticated participants.
The architecture's behavioral robustness derives substantially
from this legibility property; it asks little of the participant.

\paragraph{Paper organization.}
\Cref{sec:primitive} summarizes the settlement primitive in the
register the present paper uses.
\Cref{sec:weak-dominance} treats the weak-dominance argument's
behavioral robustness, and works out the loss-aversion analysis
at the kink.
\Cref{sec:vanhuyck} treats the Van Huyck--Battalio--Beil
coordination-failure objection and argues it does not transfer.
\Cref{sec:peer-pressure} treats the externality structure of
atomic resolution as behavioral peer pressure.
\Cref{sec:legibility} develops incentive legibility as the
architectural property that makes the design unusually
well-suited to behavioral participants.
\Cref{sec:interface} treats interface cognition: what the
architecture asks the participant to think about and how those
thinking-tasks compare to the alternatives.
\Cref{sec:empirical} states the open empirical questions.
\Cref{sec:conclusion} concludes.

% ════════════════════════════════════════════════════════════════════
\section{The Settlement Primitive}
\label{sec:primitive}

The bonded primitive has a formal mechanism-design treatment,
and its deployed Solidity implementation is verified through
TLA\textsuperscript{+}, Halmos symbolic execution, Echidna
property-based fuzzing, and Certora specification rules. We
summarize the features the present paper uses.

\paragraph{Asymmetric bonding (Mechanism~1).}
For a transaction with payment $\pay > 0$ and cumulative upstream
value $\cumval \geq \pay$, the buyer locks $2\pay$ and the seller
locks $2\cumval$. The payoff matrix in the two-party game is:

\begin{center}
\begin{tabular}{lcc}
\toprule
 & $a_B = \Coop$ & $a_B = \Def$ \\
\midrule
$a_S = \Coop$ & $(-\pay,\; +\pay)$ & $(-2\pay,\; -2\cumval)$ \\
$a_S = \Def$  & $(-2\pay,\; -2\cumval)$ & $(-2\pay,\; -2\cumval)$ \\
\bottomrule
\end{tabular}
\end{center}

Cooperation weakly dominates defection for both players.
Cooperation is the unique profile surviving iterated elimination
of weakly dominated strategies.

\paragraph{Buyer dominance with atomic resolution (Mechanism~2).}
Only the buyer can trigger resolution; resolution settles all
active orders simultaneously or not at all. The atomicity rule
induces a weakest-link subgame among sellers. Each seller's
defection imposes externality $\varepsilon_{k \to i} = \pay_i +
2\cumval_i$ on every co-seller.

\paragraph{What the behavioral analysis adds.}
The full-rationality analysis assumes participants compute the
matrix above, recognize the weak-dominance relations, and play
cooperatively. The behavioral question is whether participants
who do not perform this computation explicitly nonetheless
converge on cooperation, and how they do so. The answer, we will
argue, is that they do, through cognitive shortcuts that the
mechanism's structure makes reliable.

% ════════════════════════════════════════════════════════════════════
\section{Weak Dominance and Loss Aversion at the Kink}
\label{sec:weak-dominance}

The mechanism paper's central equilibrium result rests on weak
dominance: cooperation is at least as good as defection for both
players, with strict inequality when the counterparty cooperates.
Behavioral game theory has historically been wary of mechanisms
that rely on weak dominance, because the fully-strict case is
more robust to the cognitive limitations and non-pecuniary
preferences participants exhibit. We argue the bonded mechanism's
weak-dominance is more robust than it appears.

\paragraph{The strictness gap is at the kink, not the magnitude.}
The strict-inequality case occurs when the counterparty
cooperates: cooperation yields recovery of the bond plus the net
flow ($-\pay$ for buyer, $+\pay$ for seller); defection yields
forfeiture of the full bond ($-2\pay$ for buyer, $-2\cumval$ for
seller). The gap between these outcomes is large: for the buyer,
the gap is $\pay$; for the seller, the gap is $\pay + 2\cumval$,
which scales with chain depth via the cumulative-value
accumulator. The gap is large in absolute terms (recover the bond
or lose it entirely) and is large relative to the natural
reference point (the bond posted at commitment).

The non-strict case occurs when the counterparty defects: both
cooperation and defection yield the full bond loss because the
defection prevents resolution regardless of the participant's
play. This case is non-strict, but it is also the case in which
the participant has no decision to make: their payoff is
determined by the counterparty's defection, and their own action
is moot.

\paragraph{Loss aversion sharpens the strict case.}
Prospect theory \citep{kahneman_tversky_1979,
tversky_kahneman_1992} establishes that participants weight
losses more heavily than equivalent gains, with empirical
loss-aversion coefficients typically in the range 1.5 to 2.5.
The asymmetry is between the gain domain and the loss domain;
within each domain, prospect theory documents diminishing
sensitivity (the value function is concave over gains and convex
over losses).

The architectural question is: which feature of prospect theory
does the bonded mechanism's strict-dominance gap engage --- the
loss-aversion asymmetry, or the within-domain diminishing
sensitivity? An earlier framing of this paper conflated the two.
The correct analysis is that the bonded mechanism engages the
loss-aversion asymmetry at the kink between cooperation and
defection, while the within-domain diminishing sensitivity
operates on differences within either domain (e.g., the
difference between losing $2\pay$ now and losing $2\pay$ later)
and does not enter the architectural analysis.

The participant's choice between cooperation and defection is a
choice between a small loss (the net flow $-\pay$ for the buyer)
and a large loss (the bond forfeiture $-2\pay$). On prospect
theory's loss-domain value function, the disutility of a loss
$x \geq 0$ is approximately $\lambda \cdot |x|^{\alpha}$, where
$\lambda \in [1.5, 2.5]$ is the loss-aversion ratio (the asymmetry
between gain and loss domains at the reference point) and
$\alpha \approx 0.88$ is the loss-domain curvature parameter
\citep{tversky_kahneman_1992}. The disutility of cooperation
is $\lambda \cdot \pay^{\alpha}$ and the disutility of defection
is $\lambda \cdot (2\pay)^{\alpha} = 2^{\alpha} \cdot \lambda
\cdot \pay^{\alpha}$. The ratio of the disutilities is
$2^{\alpha} \approx 2^{0.88} \approx 1.84$, which places the
defection disutility roughly $84\%$ above the cooperation
disutility. The diminishing-sensitivity discount within the loss
domain ($\alpha < 1$) does flatten the comparison, but not
nearly enough to defeat it: $2^{\alpha} > 1$ for any
$\alpha > 0$, so the larger loss is strictly more disutile than
the smaller for any feasible curvature parameter the empirical
literature has documented. Loss aversion at the kink (the
$\lambda$ asymmetry between domains) does no work in this
comparison since both outcomes are in the loss domain; what does
the work is the magnitude difference, weighted by
$2^{\alpha} > 1$.

\paragraph{Behavioral robustness of weak dominance.}
The combined effect is that participants treat cooperation as
strictly preferred to defection in the cases where the
mechanism-paper analysis says cooperation strictly dominates,
and treat the two as equivalent in the cases where the
mechanism-paper analysis says they are equivalent (both yield
bond loss because the counterparty defected). The behavioral
prediction matches the rational-choice prediction at the
operationally-relevant choice point.

\paragraph{Non-pecuniary preferences.}
The full-rationality analysis assumes participants care only
about pecuniary payoffs. The experimental literature documents
substantial non-pecuniary preferences: fairness
\citep{fehr_schmidt_1999, bolton_ockenfels_2000}, reciprocity,
spite, and intrinsic motivation. The honest scope-claim is that
the bonded mechanism's equilibrium argument assumes pecuniary
preferences and may admit additional behavior under
non-pecuniary preferences. Three observations bound the concern.

First, fairness preferences in inequity-aversion form
\citep{fehr_schmidt_1999} predict cooperation when the symmetric
case applies (root commitment with $\pay = \cumval$, both parties
bonding $2\pay$). The bonded mechanism is symmetric at root and
asymmetric only deeper in the process tree, where the
inequity-aversion concern is the seller's higher exposure relative
to the buyer's. Inequity aversion in this direction would predict
a participant refusing to enter the bonded commitment in the first
place rather than defecting after entering; the equilibrium-of-play
result is robust to this correction at the cost of slower
multi-party-process formation.

Second, reciprocity and spite operate primarily in the
counterparty's-defection case where both cooperation and defection
yield bond loss for the participant; the participant who is
spite-motivated may prefer the defection branch as a
co-counterparty-punishment gesture. The mechanism's outcome is
unchanged (both branches yield bond loss); the participant's
internal experience may differ.

Third, intrinsic motivation to perform (a craft-based or honor-based
preference for delivering the agreed work regardless of contractual
incentive) reinforces cooperation; the bonded mechanism's pecuniary
incentive aligns with the intrinsic motivation rather than competing
against it. This is structurally favorable for the mechanism's
behavioral robustness.

% ════════════════════════════════════════════════════════════════════
\section{The Van Huyck Coordination Failure Does Not Transfer}
\label{sec:vanhuyck}

The classical experimental literature on weakest-link coordination,
beginning with \citet{vanhuyck_battalio_beil_1990}, documents a
robust empirical finding: in groups larger than two, weakest-link
coordination games converge to the Pareto-inferior all-defect
equilibrium via pessimistic beliefs about co-players. The result
appears to threaten the multi-party reading of the bonded
architecture, since the architecture is described in the mechanism
paper as inducing a weakest-link subgame among sellers.

The agent's reading of the relevant remark in the mechanism paper
flagged this potential threat directly. We develop the response
here in behavioral-game-theory register.

\paragraph{Why Van Huyck's coordination failure occurs.}
The Van Huyck--Battalio--Beil minimum-effort game asks each player
to choose an effort level from $\{1, 2, \ldots, 7\}$. Each player's
payoff is determined by the minimum effort across the group: they
gain a per-player benefit from the minimum and pay a private cost
for their own effort. In the typical experimental treatment, an
effort of 7 across the group is Pareto-optimal, but any individual
player can guarantee a non-negative payoff by choosing the minimum
effort 1 regardless of others' choices. The combination produces
multiple Nash equilibria (any common effort level), and the
empirical finding is that groups of 14--16 players converge on the
all-1 Pareto-inferior equilibrium via pessimistic beliefs about
co-players' choices.

The structural feature that produces the failure is that
\emph{individual cooperation is not dominated} by individual
defection. A player who chooses 7 against a co-player who chooses 1
takes a private cost without producing the joint benefit; choosing
7 is rational only conditional on others' choosing 7. Pessimistic
beliefs about others' rationality therefore rationalize choosing
1, and the iterated belief-formation drives the convergence.

\paragraph{Why the bonded mechanism does not exhibit it.}
The bonded mechanism's per-seller game is structurally different.
Consider the sub-game from seller $S_i$'s perspective in the
$N$-party process tree. Two cases.

\emph{Case 1: All other sellers cooperate, buyer resolves.} Then:
\begin{align*}
  u_{S_i}(\Coop) &= +\pay_i, \\
  u_{S_i}(\Def) &= -2\cumval_i.
\end{align*}
Cooperation strictly dominates defection (gap $\pay_i + 2\cumval_i$).

\emph{Case 2: At least one other seller defects, or buyer defects.}
Then:
\begin{equation*}
  u_{S_i}(\Coop) = u_{S_i}(\Def) = -2\cumval_i.
\end{equation*}
Cooperation and defection yield the same payoff because resolution
does not occur regardless of $S_i$'s action.

The combination is weak dominance: cooperation is at least as good
as defection in every co-player profile, with strict inequality in
the cooperative profile. \emph{There is no profile in which
cooperation is dominated by defection.} The pessimistic-belief
mechanism that produces Van Huyck-style coordination failure
requires individual non-dominance --- a setting in which a player
who expects others to defect prefers to defect themselves. In the
bonded mechanism, a player who expects others to defect is
indifferent between cooperating and defecting; they are not
\emph{better off} defecting. The pessimistic-belief mechanism has
nothing to operate on.

\paragraph{The architectural reading.}
The Van Huyck result documents what happens when the per-player
game is non-dominance-solvable and players reason about each
other's rationality under pessimistic priors. The bonded
mechanism is dominance-solvable at the per-player level. The
weakest-link \emph{outcome} (one defection collapses all payoffs)
is preserved; the weakest-link \emph{strategic structure} that
produces coordination failure is foreclosed. Hirshleifer's
\citep{hirshleifer1983} characterization of weakest-link
public-goods games as a structural-payoff form continues to apply;
the experimental literature's coordination-failure transfer does
not.

\paragraph{What this leaves open.}
We do not claim the bonded mechanism is fully experimentally
verified. The argument above is theoretical: the structural
condition for Van Huyck coordination failure is absent in the
bonded mechanism. Whether participants in a bonded mechanism in
fact converge on cooperation at the rates the
dominance-solvability argument predicts is an experimental
question \Cref{sec:empirical} returns to. The honest claim is
that the negative experimental finding from the existing
weakest-link literature does not transfer; the positive
experimental finding for the bonded mechanism remains to be
established empirically.

% ════════════════════════════════════════════════════════════════════
\section{Atomic Resolution as Peer Pressure}
\label{sec:peer-pressure}

The mechanism paper's analysis of atomic resolution
characterizes the externality structure formally: each seller's
defection imposes externality $\varepsilon_{k \to i} = \pay_i +
2\cumval_i$ on every other seller. The external pressure is
\emph{economic}, in the sense that it operates on each seller's
expected bond exposure. The behavioral question is whether the
pressure is also \emph{social}, in the sense that participants
experience it as peer pressure rather than only as expected-value
arithmetic. We argue it is.

\paragraph{The behavioral peer-pressure literature.}
The experimental literature on social-pressure activation
\citep{cialdini_goldstein_2004, fehr_gachter_2000} finds that
participants modify their behavior in response to peer monitoring
even at low stakes, and that the response increases with the
stakes faced by the monitor as well as by the monitored. The
seminal Fehr--G\"achter result on punishment in public-goods games
shows that participants will pay private costs to punish defectors
even when no monetary reward is available; this establishes that
peer pressure activates as a behavioral phenomenon even without an
expected-value rationale.

The bonded mechanism's externality structure activates a stronger
form of the same phenomenon. Each seller has an
\emph{expected-value} reason to monitor co-sellers: their bond is
forfeit if any co-seller defects. This is unusually high stakes:
each seller faces $\pay_i + 2\cumval_i$ in damages from any
co-seller's defection, where $\cumval_i$ scales with chain depth.
At chain depth $i = 5$ with equal payments $p$ per stage, the
damages are $11 p$, and at $i = 10$ they are $21 p$. The behavioral
peer-pressure literature operates at much lower stakes; the bonded
mechanism's stakes are an order of magnitude higher.

\paragraph{The social-substrate analogue.}
Joint-liability microfinance \citep{ghatak1999, stiglitz1990}
and the empirical literature on the Grameen Bank document a
substantial reduction in reported default rates under
joint-liability lending. \citet{grameen1999} reports a default
rate of approximately 2\% under the joint-liability regime,
compared to industry-conventional 20\% benchmarks for
unsecured rural lending; \citet{morduch_1999} contests the
headline figures on accounting grounds, arguing that loan
rescheduling and definition-of-default ambiguity make the
self-reported 2\% figure optimistic. We treat the
$20\% \to 2\%$ comparison as a stylized benchmark of the
joint-liability literature's headline finding rather than as a
peer-reviewed empirical estimate; the contemporary literature
attributes whatever reduction is real to a combination of
social peer-monitoring activation, social-sanction technology,
and reputation effects.

The mechanism paper's reduction-from-microfinance result
(Proposition~5.13 in Paper~A) shows that the bonded mechanism
reproduces the equilibrium peer-pressure outcome of joint-liability
microfinance under strictly weaker assumptions: no repeated
interaction, no local information, no exogenous punishment
technology, no joint-liability contracting. The behavioral reading
of the same result is that the social-pressure mechanism the
microfinance literature documents is, in the bonded setting,
constituted directly by the bonding rule rather than mediated
through the social substrate. The bonded mechanism is therefore not
merely \emph{equivalent} to joint-liability microfinance in
equilibrium; it produces the social-pressure dynamics through a
structurally cleaner channel that does not require the social
substrate to be in place.

\paragraph{Empirical extrapolation.}
Read as a stylized upper-bound benchmark for what
joint-liability-style peer pressure has been claimed to achieve,
the Grameen $20\% \to 2\%$ comparison gives an indication of the
order of magnitude the bonded mechanism's peer-pressure
activation might support. Three structural margins favor the
bonded mechanism over the Grameen setting: the externality
magnitude scales with chain depth (deeper sellers face larger
co-seller-defection costs); information availability is higher
(every other seller's bond exposure is on chain); and the
punishment technology is mechanical (the bonding rule rather
than social sanction). We do not claim a specific quantitative
prediction. We note that the architectural features that
produced whatever effect Grameen achieves are present, in the
bonded mechanism, in strengthened form, and that the
Morduch-style accounting concerns about Grameen's headline
figures do not apply in the same way to a mechanism whose
default and resolution events are on-chain and unambiguous.

% ════════════════════════════════════════════════════════════════════
\section{Incentive Legibility}
\label{sec:legibility}

The behavioral robustness of the bonded mechanism rests on a
property we call \emph{incentive legibility}: the cognitive
operation a participant must perform to assess the equilibrium
is a single comparison at the margin, not a multi-step iterated
reasoning about counterparty rationality.

\paragraph{The cognitive demand of standard mechanism design.}
A typical mechanism-design proof requires the participant to
reason about: their own preferences over outcomes; the
counterparty's preferences over outcomes; their belief over the
counterparty's choices; the counterparty's belief over their own
choices; the equilibrium concept the mechanism rests on; and any
common-knowledge assumptions the equilibrium requires. The
behavioral literature documents that participants do not perform
this reasoning in any reliable way --- the level-$k$ thinking
literature \citep{stahl_wilson_1995, costa-gomes_etal_2001} finds
that most participants reason at level 1 or 2 (one or two steps
of iterated thinking), not the unbounded level the
common-knowledge equilibrium concept presumes.

The architectural concern is that mechanisms whose equilibrium
properties require unbounded iterated reasoning are not robust to
the cognitive limitations of real participants. The bonded
mechanism is structurally different: the equilibrium argument
requires only level-1 reasoning.

\paragraph{The level-1 argument for cooperation.}
A participant in the bonded mechanism asks: ``If I cooperate, what
is my outcome? If I defect, what is my outcome?'' The first answer
is bounded by what the bonding rule produces conditional on the
counterparty's action; the second answer is the bond loss
unconditional on the counterparty's action. The participant
compares the two and chooses the better one. This is level-1
reasoning: one step of forward-looking analysis, no iterated
reasoning about the counterparty's reasoning, no common-knowledge
assumption.

The argument's robustness follows: a participant who can perform
level-1 reasoning --- which is the floor of strategic reasoning in
any documented adult population --- has all the cognitive
machinery the bonded mechanism requires. The mechanism does not
ask the participant to compute the iterated elimination or to
verify common knowledge of payoffs; it asks them to compare
two outcomes at the margin. The dominance argument the mechanism
paper formalizes corresponds, behaviorally, to the comparison
the participant naturally performs.

\paragraph{The Schelling-focal-point comparison.}
Schelling-focal-point coordination
\citep{schelling1960strategy, sugden1995theory} similarly works
through a level-1 argument: the participant identifies the salient
solution and chooses it, expecting the counterparty to do the
same. The bonded mechanism's coordination is even simpler than
focal-point coordination because it does not require any salience
computation. The architecturally privileged choice is determined
by the mechanism's payoff structure rather than by any framing-
or attention-dependent salience. The participant's
choice-of-cooperation is not a guess about what the counterparty
will do; it is the choice that produces the better outcome
conditional on either counterparty action.

\paragraph{Implications for the architecture's adoption.}
A mechanism's adoption in a population is partly a function of
how well the population can verify the mechanism's equilibrium
properties. A mechanism whose equilibrium argument is opaque
will struggle to attract participants regardless of the
argument's correctness, because participants who cannot verify
the argument cannot rationally trust it. The bonded mechanism's
incentive legibility is therefore not just behavioral robustness
but adoption advantage: participants can verify the argument
themselves with the cognitive operations they routinely perform.

% ════════════════════════════════════════════════════════════════════
\section{Interface Cognition}
\label{sec:interface}

The architecture's behavioral properties also depend on the
\emph{interface} through which participants interact with the
mechanism. A clean equilibrium argument can be defeated by an
interface that obscures the relevant choice or surfaces irrelevant
information. We treat interface cognition briefly because it is
where the architecture meets real participants.

\paragraph{What the architecture asks the participant to think about.}
The participant in a bonded commitment is asked to attend to four
things at the moment of commitment: (i) the goods or services
being exchanged, (ii) the price, (iii) the bond posture (which is
mechanically determined by the price and the cumulative-value
context), and (iv) the resolution path (which is the buyer's
signature plus any clauses the parties have agreed). The mechanism
paper's invariants are encoded in the structure rather than left
to the participant's verification: the bonding rule produces
$2\pay$ and $2\cumval$ automatically, the buyer-dominance rule
constrains who can resolve, and the atomic-resolution rule
constrains how resolution proceeds.

The participant does not need to verify the mechanism; the
mechanism enforces itself. The participant's cognitive task is
limited to the substantive commercial decision (is this exchange
worth this price?) plus the bond-feasibility decision (is the bond
within my treasury capacity?). Both are routine commercial
decisions that participants are well-equipped to perform.

\paragraph{Comparison with the platform-mediated alternative.}
A participant in a contemporary platform-mediated commerce
arrangement is typically asked to attend to: the goods or
services, the price, the platform's terms-of-service, the
platform's complaint-resolution process, the platform's
recommendation-signal mechanics, the platform's reputation system,
the platform's payment-flow timing, and the platform's
jurisdictional-coverage rules. Most of these are not substantive
commercial decisions; they are infrastructural-property
verifications the platform asks the participant to perform on its
behalf. The cognitive load is correspondingly higher.

The bonded architecture's interface cognition is structurally
lower because the architectural properties are enforced rather
than verified. A participant does not need to read the
terms-of-service of a permissionless protocol; the protocol's
behavior is its terms-of-service, and the protocol cannot violate
its terms because the terms are encoded in the bytecode. This
shifts the participant's cognitive load from infrastructural
verification to substantive commercial decision-making, which is
the right direction for a coordination architecture.

\paragraph{The risk of over-loading the interface.}
A counter-argument is that the bonded architecture's interface
may over-load participants who are unfamiliar with cryptographic
key management, on-chain transactions, gas estimation, and
similar substrate-tier concerns. This is a real concern at the
runtime tier and is the work of the wallet, signing, and onboarding
infrastructure rather than of the kernel itself. The kernel is
incentive-legible; the runtime tier may or may not be, depending
on how it presents the kernel to participants. Interface design
at the runtime tier is therefore where the architecture's
behavioral promises are kept or broken in practice. The kernel
makes the promises tractable; the runtime delivers them.

% ════════════════════════════════════════════════════════════════════
\section{Open Empirical Questions}
\label{sec:empirical}

The architecture's behavioral analysis above is theoretical. The
experimental questions it invites are several.

\paragraph{Bilateral cooperation rates.}
The first-order empirical question is whether participants in a
bonded bilateral commitment in fact converge on cooperation at
rates significantly higher than the conventional contractual
baseline. The experimental design is a controlled bilateral
exchange where one treatment uses bonded commitment and the
control uses conventional contract; the dependent variable is
the cooperation rate. We predict significantly higher cooperation
under the bonded treatment, with the loss-aversion-amplified kink
producing the bulk of the effect.

\paragraph{Multi-party coordination success.}
The second-order question is whether $N$-party processes resolve
successfully at rates that match the dominance-solvability
prediction. The Van Huyck literature establishes that conventional
multi-party coordination decays with $N$; the present paper's
argument is that the bonded mechanism does not exhibit this decay
because the per-player game is dominance-solvable. An experiment
varying $N$ across treatments would test the prediction directly.

\paragraph{Peer-pressure activation.}
The third-order question is whether the externality structure of
atomic resolution is behaviorally activated as peer pressure. The
experimental design measures whether participants in an $N$-party
bonded process modify their behavior in response to other
participants' visible bond exposure --- whether peer monitoring
activates at the rates the externality magnitude would predict.
The Grameen literature provides the upper-bound benchmark; we
expect the bonded mechanism's activation to fall between the
Grameen rate and the conventional-contract rate, with the
position depending on how much the social-substrate channel
contributed to Grameen relative to the joint-liability channel.

\paragraph{Cognitive load measurement.}
The fourth-order question is the architecture's interface
cognition relative to the alternatives. The experimental design
measures cognitive load (response time, error rate,
self-reported difficulty) across treatments where participants
interact with bonded commerce, platform-mediated commerce, and
direct off-platform commerce. We predict the bonded treatment's
load is intermediate between platform-mediated (highest, due to
infrastructural verification) and direct off-platform (lowest,
due to absence of architectural superstructure). Whether the
architecture's load is significantly lower than platform-mediated
in operational deployment is the question that determines the
adoption story.

\paragraph{Field measurements.}
Beyond the laboratory, field measurements of the architecture's
operational performance in deployments (the Spirit Air assembly
of Paper~J, the TradeLens replacement of Paper~K, the
local-commerce reference assembly of the runtime documentation)
will produce the binding evidence on whether the behavioral
predictions hold in operational settings. Laboratory results are
suggestive; field results are dispositive.

% ════════════════════════════════════════════════════════════════════
\section{Conclusion}
\label{sec:conclusion}

The bonded settlement primitive's full-rationality equilibrium
analysis carries over to the behavioral setting through four
specific properties. The weak-dominance argument operates at the
kink between cooperation and defection rather than at the
magnitude differences within either domain, where loss aversion
sharpens the comparison rather than blunting it. The Van Huyck
coordination-failure result does not transfer because the
per-player game is dominance-solvable. The atomic-resolution
externality structure activates the same peer-pressure dynamics
the Grameen literature documents in joint-liability microfinance,
through a cleaner architectural channel that does not require the
social substrate to be in place. The architecture exhibits
incentive legibility: the cognitive operation the mechanism asks
of participants is a single comparison at the margin, well within
the documented capabilities of unsophisticated participants.

The four properties together produce a mechanism whose behavioral
robustness derives from its structural simplicity. The mechanism
is not optimal in any technical sense; it is robust in the
behavioral sense that it does not ask participants to do
cognitive work they cannot reliably do. This is a different
property than rational-choice optimality, and it is, we suggest,
the more important property for any mechanism that aspires to
operate at population scale among real participants.

The contribution of this paper is the behavioral analysis.
Combined with the mechanism-design treatment, the argument is
that the bonded primitive is both formally elegant under full
rationality and behaviorally robust
under documented departures from full rationality. Whether the
behavioral predictions hold experimentally and in field deployment
is the question the architecture invites and that we identify as
the experimental work future research should undertake.

\section*{Acknowledgments}
The behavioral analysis builds on the mechanism-design work and
on the externality-engagement treatments in the political-economy
and coercion literatures.
I thank the experimental-economics research community whose work
on weakest-link coordination, loss aversion, and peer pressure
provides the empirical anchor against which the present analysis
is calibrated.

% ════════════════════════════════════════════════════════════════════
\begin{thebibliography}{99}

\bibitem[Camerer(2003)]{camerer2003behavioral}
Camerer, C.~F.
\newblock \emph{Behavioral Game Theory: Experiments in Strategic
  Interaction}.
\newblock Princeton University Press, Princeton, NJ, 2003.

\bibitem[Kahneman and Tversky(1979)]{kahneman_tversky_1979}
Kahneman, D. and Tversky, A.
\newblock Prospect Theory: An Analysis of Decision under Risk.
\newblock \emph{Econometrica}, 47(2):263--291, 1979.

\bibitem[Tversky and Kahneman(1992)]{tversky_kahneman_1992}
Tversky, A. and Kahneman, D.
\newblock Advances in Prospect Theory: Cumulative Representation
  of Uncertainty.
\newblock \emph{Journal of Risk and Uncertainty}, 5(4):297--323, 1992.

\bibitem[Fehr and Schmidt(1999)]{fehr_schmidt_1999}
Fehr, E. and Schmidt, K.~M.
\newblock A Theory of Fairness, Competition, and Cooperation.
\newblock \emph{Quarterly Journal of Economics}, 114(3):817--868, 1999.

\bibitem[Bolton and Ockenfels(2000)]{bolton_ockenfels_2000}
Bolton, G.~E. and Ockenfels, A.
\newblock ERC: A Theory of Equity, Reciprocity, and Competition.
\newblock \emph{American Economic Review}, 90(1):166--193, 2000.

\bibitem[Van Huyck et al.(1990)]{vanhuyck_battalio_beil_1990}
Van Huyck, J.~B., Battalio, R.~C., and Beil, R.~O.
\newblock Tacit Coordination Games, Strategic Uncertainty, and
  Coordination Failure.
\newblock \emph{American Economic Review}, 80(1):234--248, 1990.

\bibitem[Hirshleifer(1983)]{hirshleifer1983}
Hirshleifer, J.
\newblock From Weakest-Link to Best-Shot: The Voluntary Provision
  of Public Goods.
\newblock \emph{Public Choice}, 41(3):371--386, 1983.

\bibitem[Cialdini and Goldstein(2004)]{cialdini_goldstein_2004}
Cialdini, R.~B. and Goldstein, N.~J.
\newblock Social Influence: Compliance and Conformity.
\newblock \emph{Annual Review of Psychology}, 55:591--621, 2004.

\bibitem[Fehr and G\"achter(2000)]{fehr_gachter_2000}
Fehr, E. and G\"achter, S.
\newblock Cooperation and Punishment in Public Goods Experiments.
\newblock \emph{American Economic Review}, 90(4):980--994, 2000.

\bibitem[Ghatak(1999)]{ghatak1999}
Ghatak, M.
\newblock Group Lending, Local Information and Peer Selection.
\newblock \emph{Journal of Development Economics}, 60(1):27--50, 1999.

\bibitem[Stiglitz(1990)]{stiglitz1990}
Stiglitz, J.~E.
\newblock Peer Monitoring and Credit Markets.
\newblock \emph{World Bank Economic Review}, 4(3):351--366, 1990.

\bibitem[Yunus(1999)]{grameen1999}
Yunus, M.
\newblock \emph{Banker to the Poor: Micro-Lending and the Battle
  Against World Poverty}.
\newblock PublicAffairs, New York, 1999.

\bibitem[Morduch(1999)]{morduch_1999}
Morduch, J.
\newblock The Microfinance Promise.
\newblock \emph{Journal of Economic Literature}, 37(4):1569--1614, 1999.

\bibitem[Stahl and Wilson(1995)]{stahl_wilson_1995}
Stahl, D.~O. and Wilson, P.~W.
\newblock On Players' Models of Other Players: Theory and Experimental
  Evidence.
\newblock \emph{Games and Economic Behavior}, 10(1):218--254, 1995.

\bibitem[Costa-Gomes et al.(2001)]{costa-gomes_etal_2001}
Costa-Gomes, M., Crawford, V.~P., and Broseta, B.
\newblock Cognition and Behavior in Normal-Form Games: An Experimental
  Study.
\newblock \emph{Econometrica}, 69(5):1193--1235, 2001.

\bibitem[Schelling(1960)]{schelling1960strategy}
Schelling, T.~C.
\newblock \emph{The Strategy of Conflict}.
\newblock Harvard University Press, Cambridge, MA, 1960.

\bibitem[Sugden(1995)]{sugden1995theory}
Sugden, R.
\newblock A Theory of Focal Points.
\newblock \emph{Economic Journal}, 105(430):533--550, 1995.

\end{thebibliography}

\end{document}
